\chapter{插入、更新和删除}
\section{插入默认值}
在值列表中，关键字 DEFAULT 会在相应列中插入创建表时指定的默认值。所有 DBMS 都支持这个关键字。

如果给表中的所有列都指定了默认值（就像前面的 D 表那样），那么 MySQL、PostgreSQL 和 SQL Server 用户还有另一种选择，就是使用空的 VALUES 列表（用于 MySQL）或 DEFAULT VALUES 子句（用于 PostgreSQL 和 SQL Server）来创建一个全为默认值的新行。否则，就需要使用 DEFAULT 分别将各列设置为默认值。

如果表中既有默认值列又有非默认值列，那么也可以轻松地将列值设置为默认值。为此，只需在插入列表中不指定该列，这样就无须使用关键字 DEFAULT 了。
\section{将一张表中的行复制到另一张表中}
在 INSERT 语句中将生成所需行的查询用作值列表。

只需在 INSERT 语句后面跟一个返回所需行的查询。如果要复制源表中所有的行，那么可以在查询中不指定 WHERE 子句。与常规插入一样，并非必须显式地指定要插入哪些列，但如果没有指定目标列，则必须在所有列中都插入数据，还必须注意 SELECT 列表中值的排列顺序。

\section{复制表定义}
\begin{description}
    \item[DB2] 在 CREATE TABLE 命令中使用 LIKE 子句。DB2 的 CREATE TABLE...LIKE 命令让你能够轻松地以一张表为模板创建另一张表。为此，在关键字 LIKE 后面指定模板表的名称即可。
    \item[Oracle、MySQL 和 PostgreSQL] 在 CREATE TABLE 命令中使用一个不返回任何行的子查询。使用 Create Table As Select（CTAS）时，查询返回的行都将用来填充新建的表，除非你在 WHERE 子句中指定了一个无法满足的条件。在解决方案中，WHREE 子句中的表达式 1 = 0 导致查询不会返回任何行。因此，这条 CTAS 语句的结果是一张空表，其包含的列与 SELECT 子句指定的列相同。
    \item[SQL Server] 在一个不返回任何行的子查询中使用 INTO 子句。使用 INTO 复制表时，查询返回的行都将用来填充新建的表，除非你在 WHERE 子句中指定了一个无法满足的条件。在解决方案中，谓词中的表达式 1 = 0 导致查询不会返回任何行。因此结果是一张空表，其包含的列与 SELECT 子句指定的列相同。
\end{description}
\section{同时插入多张表}
\begin{description}
    \item[Oracle] 使用 INSERT ALL 语句或 INSERT FIRST 语句。这两条语句的语法相同，唯一的差别是一条语句使用的是关键字 ALL，另一条语句使用的是关键字 FIRST。Oracle 多表插入使用 WHEN-THEN-ELSE 子句来检查嵌套的 SELECT 返回的行，以便将它们插入正确的表中。INSERT ALL 和 INSERT FIRST 的效果几乎相同，但它们之间存在一个不同之处。INSERT FIRST 在遇到满足的条件后，就结束 WHEN-THEN-ELSE 检查，而 INSERT ALL 会检查所有的条件，即便已经遇到了满足的条件。因此，可以使用INSERT ALL 将相同的行插入多张表中。
    \item[DB2] 使用 UNION ALL 合并所有的目标表，以创建一个内嵌视图，再将数据插入这个内嵌视图中。还必须给目标表设置约束条件，确保每行都被插入正确的目标表中。DB2 解决方案有点儿不合常规。它要求给目标表定义约束条件，以确保从子查询返回的每一行都进入正确的表中。这里的诀窍是使用 UNION ALL 合并所有目标表，以生成一个视图，然后将行插入这个视图中。如果给目标表指定的约束条件不是唯一的（比如有多张表的约束条件相同），那么 INSERT 语句将不知道该将行插入哪张表中，进而以失败告终。
\end{description}
\section{禁止在特定列中插入值}
创建一个基于表的视图，只暴露那些你想暴露的列，然后规定只能通过这个视图来插入。授予用户和程序访问这个视图的权限，让它们能够填充该视图中的列。不授予用户在原表中插入行的权限。

当在简单视图中插入行时，数据库服务器将把它转换为针对底层表的插入操作。视图插入是一个复杂的主题，除非视图很简单，否则规则将很快变得非常复杂。如果你要使用视图插入功能，那么务必参阅相关的数据库文档，并确保已完全理解。
\section{修改表中的记录}
使用 UPDATE 语句来修改数据库表中既有的行。在 UPDATE 语句中，使用 WHERE 子句来指定要更新哪些行。如果没有指定 WHERE 子句，则将更新所有行。进行大规模更新前，你可能想预览结果。为此，可以执行一
条 SELECT 语句，并在其中使用你打算在 SET 子句中使用的表达式。